\uuid{0RlI}
\titre{Dénombrement}
\auteur{bodin ; hanani ; mzari}
\organisation{exo7}
\licence{CC BY-NC-SA 4.0}
\datecreate{2020-01-20}
\niveau{L1}
\module{Algèbre}
\chapitre{Logique, ensemble, raisonnement}
\sousChapitre{Ensemble}
\difficulte{3}
\theme{Ensembles, applications, dénombrement}
\exo7qcmid{lille-1.ensembles.027}

\contenu{
\begin{question}
\qclassification{100.02,101.01,102.01,102.02}
\qtags{motcle=dénombrement}
Soit $E$ un ensemble à $n$ éléments et $a\in E$. On note $\mathcal{P}_a(E)$ l'ensemble des parties de $E$ qui contiennent $a$. Quel est le cardinal de $\mathcal{P}_a(E)$ ?
\begin{answers}
    \bad{$\mathrm{Card}(\mathcal{P}_a(E))=n-1$}
    \bad{$\mathrm{Card}(\mathcal{P}_a(E))=n$}
    \good{$\mathrm{Card}(\mathcal{P}_a(E))=2^{n-1}$}
    \bad{$\mathrm{Card}(\mathcal{P}_a(E))=2^n$}
\end{answers}
\begin{explanations}
Les éléments de $\mathcal{P}_a(E)$ sont de la forme $\{a\}\cup A$ où $A\subset E\setminus \{a\}$. Donc $\mathrm{Card}(\mathcal{P}_a(E))=\mathrm{Card}(\mathcal{P}(E\setminus \{a\}))=2^{n-1}$.
\end{explanations}
\end{question}
}
